\begin{englishabstract}
Brain-computer interface (BCI) IoT links carry highly sensitive neural data and control commands, and are exposed to replay, tampering, man-in-the-middle (MITM), and traffic-fingerprinting attacks. While TLS/DTLS-, MQTT/TLS-, and CoAP/DTLS-based deployments are cryptographically mature, constrained devices and restricted public networks still face engineering trade-offs among resource cost, reachability, and traffic observability. This thesis investigates a combinatorial Sudoku protocol prototype and evaluates a combined design consisting of authenticated handshake, AEAD record protection, and a Sudoku-based reversible traffic-appearance layer with randomized padding.

The protocol is implemented as an enhancement layer between the application and TCP. The handshake integrates certificate-based authentication, ephemeral key agreement, transcript binding, and confirmation MAC; the session layer adopts AEAD records; the outer appearance layer provides reversible byte mapping, randomized padding, an ASCII mode, and customizable templates. The evaluation pipeline includes local baseline benchmarking, fingerprinting and active-attack validation, and cross-region restricted-network experiments under unified metrics and repeated-run statistics.

The evaluation shows that the proposed scheme maintains latency at the same practical level as mainstream secure protocols under restricted-network conditions, while exhibiting lower data-plane memory usage. In security validation, replay and tampering inputs are consistently rejected. In side-channel evaluation, MQTT/TLS, CoAP/UDP, Pure-AEAD, and the unpadded Sudoku mode all show 100\% class recoverability, whereas randomized padding in the pure Sudoku mode reduces the recovery rate to 16.67\%. The comparison further indicates complementary engineering trade-offs between the SP mode (obfuscation-oriented) and the SK mode (bandwidth-oriented).

This thesis contributes an end-to-end chain from mathematical modeling to prototype implementation and experimental validation, and explicitly reports both strengths and limitations against mainstream baselines. The results provide a practical reference for protocol selection and future standardization in secure BCI-IoT communication.

\englishkeyword{BCI, IoT secure communication, Sudoku coding, random padding, traffic fingerprinting}
\end{englishabstract}
